\documentclass[12pt]{article}
\usepackage{titling}
\newcommand{\subtitle}[1]{%
  \posttitle{%
    \par\end{center}
    \begin{center}\large#1\end{center}
    \vskip0.5em}%
}
\usepackage[margin=1in]{geometry}
\usepackage{amsmath,amssymb,amsthm,amsfonts}
\usepackage{graphicx}
\usepackage{hyperref}
\usepackage{framed}
\hypersetup{colorlinks=true, linkcolor=blue, urlcolor=blue, citecolor=blue}

% Chapter-specific environments
\theoremstyle{definition}
\newtheorem{definition}{Definition}[section]
\newtheorem{remark}{Remark}[section]
\newtheorem{theorem}{Theorem}[section]
\newtheorem{lemma}[theorem]{Lemma}
\newtheorem{corollary}[theorem]{Corollary}

\newenvironment{editorsnote}{%
  \begin{framed}
  \noindent\textbf{Editor's Note:}\em
}{\end{framed}}

\title{\textbf{Chapter 2: An Initial, Concrete Exploration}}
\author{Dmytro Surdu}
\date{}

\begin{document}
\maketitle

\begin{editorsnote}
This chapter presents a verbatim reprint of an early research paper that initiated our inquiry. Its goal was to explore the boundary between tractable and intractable problems in a statistical setting. The paper focuses on a very concrete scenario: what can be said about the stability of a sample's mean and standard deviation when new data points are known to be "clipped," or confined to a band around the current mean?

The main result (Theorem 6.2) is that if the \emph{entire} future data stream is given as input, the problem of checking its conformity is computationally trivial (in $\mathbf{P}$).

Pay close attention to the limitations discussed in the paper's conclusion and in our annotations. This work successfully simplifies the "universal quantifier" problem by assuming the future is fully known. This cleanly demonstrates the tractability of verification in a world without uncertainty, but it leaves the core metrological question unanswered: **How can we certify such a property when the future is unknown and we only have a finite witness?** This paper's limitations are the direct motivation for the more powerful tools we will develop in the subsequent chapters.
\end{editorsnote}

\hrule
\vspace{2em}


\subtitle{Tractability, Testability, and Tail Restrictions:\\ A Formal Analysis of Certificate--Based Verification Problems}
\author{Anonymous ChatGPT Collaboration}
\date{Original Draft Version}


\maketitle

\begin{abstract}
We synthesise a sequence of results developed interactively in conversation, examining (i) the conceptual distinction between tractability and testability in computational complexity, (ii) the difficulty of certifying distributional hypotheses under tail--element restrictions, (iii) the position of such languages in the arithmetic and polynomial hierarchies, and (iv) a precise algebraic analysis of how small ``clipped'' additions to a sample affect its mean and unbiased standard deviation.  We culminate by classifying the natural decision problem obtained from these statistics as lying in the class $\mathbf{P}$ (hence in $\mathbf{NP}$ but not $\mathbf{NP}$--complete unless $\mathbf{P} = \mathbf{NP}$).
\end{abstract}

\tableofcontents

\section{Introduction}
The perennial question ``Can we verify a claim more easily than we can find a proof of it?'' permeates both theoretical computer science and applied statistics.  In complexity theory this dichotomy manifests as the classes $\mathbf{P}$ (tractable, efficiently solvable) and $\mathbf{NP}$ (efficiently verifiable once a short witness is supplied).  Statistics faces an analogous gap between determining whether data plausibly originate from a hypothesised distribution and confirming with confidence that a given sample obeys certain distributional properties.  

In a series of informal conversations we investigated the interaction of these themes, guided by the following overarching questions:
\begin{enumerate}
  \item What constitutes an $\mathbf{NP}$ witness for the statement ``the datum $x$ could have been generated by distribution $q$''?
  \item When do tail restrictions on future samples suffice to collapse a universally quantified property into $\mathbf{NP}$ (or even $\mathbf{P}$)?
  \item How do elementary clipping rules affect fundamental sample statistics, and what complexity class contains the resulting decision problems?
\end{enumerate}

This paper formalises and extends those results.

\section{Preliminaries}
\subsection{Complexity Classes}
We recall only those definitions required for later arguments; see \cite{sipser2012introduction} for a comprehensive treatment.

\begin{definition}[Polynomial--time decidability]
A language $L$ is in $\mathbf{P}$ if there exists a deterministic Turing machine $M$ and a polynomial $p$ such that for all $x$, $M$ halts on input $x$ in time $\le p(|x|)$ and accepts iff $x \in L$.
\end{definition}

\begin{definition}[Verifiability -- $\mathbf{NP}$]
A language $L$ is in $\mathbf{NP}$ if there exist a polynomial $p$ and a deterministic verifier $V$ such that $x \in L$ iff $\exists w \in \{0,1\}^{\le p(|x|)}$ with $V(x,w)=\mathrm{accept}$, and $V$ runs in time $\le p(|x|)$.
\end{definition}

\begin{definition}[Recursively Enumerable (r.e.)]
A set $S \subseteq \Sigma^*$ is \emph{recursively enumerable} if there exists a Turing machine that enumerates exactly the elements of $S$ (possibly without halting).  Equivalently, $S$ is $\Sigma^0_1$ in the arith­metical hierarchy.
\end{definition}

Throughout we adopt the usual inclusions $\mathbf{P} \subseteq \mathbf{NP} \subseteq \Sigma^0_1$.

\subsection{Statistical Notation}
For a finite sample $S_m = (s_1,\dots,s_m)$ with $m\ge 2$ we write
\[
  \mu_m = \frac1m \sum_{k=1}^m s_k,\qquad
  \sigma_m^2 = \frac1{m-1}\sum_{k=1}^m (s_k-\mu_m)^2 .
\]
The unbiased sample standard deviation is $\sigma_m = \sqrt{\sigma_m^2}$.

\section{Certificates for Distributional Support}
\label{sec:certificates}
Consider the language
\[
  L_{\mathrm{support}} = \{(x,q)\mid q(x)>0\}.
\]
If $q$ is specified by a probabilistic polynomial--time sampler \texttt{Sample}$_q(r)$, one obtains an $\mathbf{NP}$ witness simply by revealing the random tape $r$ that satisfies \texttt{Sample}$_q(r)=x$.  Verification reruns the sampler deterministically; see \cite[Sec.~2]{goldreich2008foundations}.

\section{Tail Restrictions and the Arithmetic Hierarchy}
Let a fixed per--item predicate $P$ define an admissible ``clipping'' band, e.g.
\[
  P(y) : |y-\mu_q| < 3\sigma_q .
\]
Given a finite prefix $u \in \Sigma^*$, write
\[
  u \in L_P \;\Longleftrightarrow\; \forall v \in \Sigma^{\omega}: P\bigl((u\cdot v)[i]\bigr) \text{ for all } i\ge |u|.
\]
Because of the universal quantifier over an unbounded tail, $L_P$ is a $\Pi^0_2$ set and therefore \emph{not} recursively enumerable.  Its complement --- ``there exists a violation of $P$'' --- is $\Sigma^0_1$ and hence r.e.

\section{Effect of Clipping on Sample Statistics}
We now turn to the algebraic core.
\subsection{One--Step Update}
\begin{theorem}[Single Clipped Addition]
\label{thm:one_step}
Let $a\ge 1$ and assume a new observation $s_{n+1}$ obeys $|s_{n+1}-\mu_n|<a\sigma_n$.  Then
\begin{equation}\label{eq:mean_bound}
  |\mu_{n+1}-\mu_n| < a\sigma_n,
\end{equation}
\begin{equation}\label{eq:sd_bound}
  \sigma_{n+1} < a\sigma_n .
\end{equation}
\end{theorem}
\begin{proof}
Equation~\eqref{eq:mean_bound} follows immediately from
$\mu_{n+1}=\mu_n+\dfrac{s_{n+1}-\mu_n}{n+1}$.  For the variance we have
\[
  \sigma_{n+1}^2 = \left(1-\frac1n\right)\sigma_n^2 + \frac{n}{(n+1)^2}(s_{n+1}-\mu_n)^2,
\]
which, under $|s_{n+1}-\mu_n|<a\sigma_n$, implies
$\sigma_{n+1}^2<\sigma_n^2\bigl(1-\tfrac1n+a^2\tfrac{n}{(n+1)^2}\bigr)\le a^2\sigma_n^2$ whenever $a\ge 1$.
\end{proof}

\subsection{Two--Step and Inductive Extension}
\begin{theorem}[Stability Under Repeated Clipping]
\label{thm:two_step}
Let $a\ge 1$ and suppose $s_{n+1},s_{n+2}$ each satisfy the clip rule
$|s_{k}-\mu_n|<a\sigma_n$.  Then the enlarged sample $S_{n+2}$ obeys
\eqref{eq:mean_bound} and \eqref{eq:sd_bound} with $(n+1)$ replaced by $(n+2)$.  Consequently, by induction, \emph{every} finite extension assembled via rule $P$ preserves the bounds.
\end{theorem}
\begin{proof}
Set $\delta_1=s_{n+1}-\mu_n$ and $\delta_2=s_{n+2}-\mu_n$.  The new mean satisfies
$|\mu_{n+2}-\mu_n|\le\frac{|\delta_1|+|\delta_2|}{n+2}<\frac{2a\sigma_n}{n+2}<a\sigma_n$ for $n\ge2$.

For the variance, writing $\mathrm{SSE}_m=(m-1)\sigma_m^2$ gives
\[
\sigma_{n+2}^2\le\frac{(n-1)\sigma_n^2+\delta_1^2+\delta_2^2}{n+1}
<\frac{n-1+2a^2}{n+1}\sigma_n^2\le a^2\sigma_n^2 \quad(a\ge1).
\]
\end{proof}

\section{Complexity of the Statistical Decision Problem}
\begin{definition}[Clipped--Sample Verification Language]
Fix $a>1$ and $n\ge 2$.  Let an instance be $(S_n,S_i)$ with $i\ge n$.  Define
\[
  L_{\mathrm{clip}} = \bigl\{(S_n,S_i)\mid |\mu_i-\mu_n|<a\sigma_n \text{ and } \sigma_i < a\sigma_n\bigr\}.
\]
\end{definition}

\begin{theorem}
$L_{\mathrm{clip}} \in \mathbf{P}$.  Hence $L_{\mathrm{clip}} \in \mathbf{NP}$ but is not $\mathbf{NP}$--complete unless $\mathbf{P} = \mathbf{NP}$.
\end{theorem}
\begin{proof}

\subsection*{1. Input Model}

Let
\[
S_i = (s_1,\dots,s_i)
\]
be an input instance, where each measurement
\[
s_k = \frac{p_k}{q_k}
\]
is a rational in lowest terms.  Assume each numerator \(p_k\) and denominator \(q_k\) has bit‐length at most \(B\).  Then the total input size is
\[
N = i \times O(B).
\]

\subsection*{2. Running‐Sum Representation}

We maintain two pairs of integers:
\[
\begin{aligned}
A_n &= \sum_{k=1}^n p_k \,\Bigl(\prod_{j\neq k} q_j\Bigr),
&\quad
D_n &= \prod_{k=1}^n q_k,\\
S_n &= \sum_{k=1}^n \bigl(p_k\prod_{j\neq k} q_j\bigr)^2,
&\quad
D_n^2 &= \bigl(\!D_n\bigr)^2.
\end{aligned}
\]
Then
\[
\sum_{k=1}^n s_k = \frac{A_n}{D_n},
\quad
\sum_{k=1}^n s_k^2 = \frac{S_n}{D_n^2}.
\]
We update these in each step:

\[
\begin{aligned}
A_{n+1} &= A_n\cdot q_{n+1} + p_{n+1}\cdot D_n,\\
D_{n+1} &= D_n\cdot q_{n+1},
\end{aligned}
\qquad
\begin{aligned}
S_{n+1} &= S_n\cdot q_{n+1}^2 + \bigl(p_{n+1}\prod_{j=1}^{n}q_j\bigr)^2,\\
D_{n+1}^2 &= \bigl(D_{n+1}\bigr)^2.
\end{aligned}
\]

\subsection*{3. Bit‐Length Growth and Operation Cost}

After \(n\) steps:
\[
\log_2 D_n = O(nB),
\quad
\text{and}\quad
\text{bit‐length}(A_n),\,\text{bit‐length}(S_n) = O(nB).
\]
Multiplication or addition of two \(L\)-bit integers on a Turing machine costs \(O(L^2)\) time.  Thus each update
\[
(A_n,D_n)\to(A_{n+1},D_{n+1})
\quad\text{and}\quad
(S_n,D_n^2)\to(S_{n+1},D_{n+1}^2)
\]
takes 
\[
O\bigl((nB)^2\bigr)
\]
time.  Summing over \(n=1\) to \(i\) gives
\[
\sum_{n=1}^i O\bigl((nB)^2\bigr)
= O\bigl(i^3B^2\bigr).
\]
Since \(N = iB\), this is \(O(N^3)\).

\subsection*{4. Computing Mean and Variance}

We obtain exact rationals:
\[
\mu_i = \frac{A_i}{i\,D_i},
\quad
\sigma_i^2 = \frac{i\,S_i - A_i^2/D_i}{\,i(i-1)\,D_i^2}.
\]
Each involves a fixed number of multiplications, additions, and comparisons on \(O(iB)\)-bit integers, costing \(O((iB)^2)=O(N^2)\) time.

\subsection*{5. Verifying the Clipped‐Sample Inequalities}

To decide membership in
\[
L_{\mathrm{clip}} = \{(S_n,S_i)\mid |\mu_i - \mu_n| < a\,\sigma_n,\;\sigma_i < a\,\sigma_n\},
\]
we compute \(\mu_n,\sigma_n\) and \(\mu_i,\sigma_i\) as above, then perform two integer comparisons.  Total time remains polynomial, specifically
\[
O(N^3) + O(N^2) = O(N^3).
\]

\subsection*{6. Conclusion}

All steps—from updating running sums to testing the two inequalities—run in time polynomial in the input length \(N\).  Hence
\[
L_{\mathrm{clip}} \in \mathbf{P}.
\]

\end{proof}

\section{Discussion and Future Work}
The investigation highlights a recurring moral: \emph{universal quantification over unseen data lifts decision problems out of $\mathbf{NP}$ unless one supplies a finite artefact that fixes the data}.  In our setting the per--element clipping rule $P$ constrains but does not determine the tail, so the language of ``all tails good'' lies in $\Pi^0_2$.  Conversely, once the tail is given explicitly (as in $L_{\mathrm{clip}}$) the problem collapses to $\mathbf P$.

An immediate extension is to study \emph{interactive} or \emph{zero--knowledge} variants: can a prover convince a verifier that an \emph{infinite} stream will forever obey $P$ while revealing only $O(1)$ information?  Such protocols inhabit the classes $\mathbf{AM}$, $\mathbf{IP}$, or $\mathbf{SZK}$ and merit further exploration.

\section*{Acknowledgements}
The authors thank the conversational partner for insisting on rigorous algebraic detail.

\bibliographystyle{plain}
\begin{thebibliography}{9}
\bibitem{sipser2012introduction}
M.~Sipser.
\newblock \emph{Introduction to the Theory of Computation}, 3rd~ed.
\newblock Cengage, 2012.

\bibitem{goldreich2008foundations}
O.~Goldreich.
\newblock \emph{Foundations of Cryptography, Vol.~2}.
\newblock Cambridge University Press, 2008.
\end{thebibliography}

\end{document}

% --- End of Verbatim Reprint ---
\vspace{2em}
\hrule

\begin{editorsnote}
\textbf{Concluding Analysis.} The central argument of this paper is sound but limited. It correctly identifies the complexity of statements with universal quantifiers ($\Pi^0_2$) and then shows that if this quantification is removed—by providing the full future sample $S_i$ as part of the input—the problem becomes tractable ($\mathbf{P}$).

This highlights the core challenge for a theory of metrological certification:
\begin{itemize}
    \item The witness in Theorem 6.2 is the data stream $S_i$ itself. This is not a "short" witness in the sense of $\mathbf{NP}$. If the future stream is long, the witness is long.
    \item The true goal is to find a witness that is \emph{logarithmically} or \emph{polynomially} small relative to the length of the future stream it purports to certify.
\end{itemize}
This paper shows what happens when we solve the problem by fiat. The next chapter will introduce the information-theoretic tools necessary to solve it in a more meaningful way, by finding a truly short witness that can substitute for the entire, infinite future.
\end{editorsnote}

\end{document}